\documentclass[a4paper, 11pt]{article} %Here you can change the font size of the entire document
%\\------------------------------------------------------\\
%\\ Instructions to user
% Please add the following build script into your LaTeX compiler
%\\ In Texmaker navigate to Options->Configure Texmkaer->Quick Build and paste this script into the "User" section
%\\ pdflatex -synctex=1 -interaction=nonstopmode %.tex|pdflatex -synctex=1 -interaction=nonstopmode %.tex|pdflatex -synctex=1 -interaction=nonstopmode %.tex|makeindex %.nlo -s nomencl.ist -o %.nls|pdflatex -synctex=1 -interaction=nonstopmode %.tex|"C:/Program Files/Adobe/Reader 11.0/Reader/AcroRd32.exe" %.pdf


%\\------------------------------------------------------\\
%\\ Packages needed for the operation of the template
%\\ DO NOT DELETE THIS
\usepackage{sectsty} %Package needed to format the section headings. This package is responsible for the \underbar and \underline errors
\usepackage{graphicx} %Package for importing graphics (pics)
\usepackage{ragged2e} %Package for justifying text
\usepackage{tikz} %Package for advanced image positioning
\usepackage{xcolor} %Package for changing the color of text
\usepackage{etoolbox} %Package to change abstract title to be at the top of the page
\patchcmd{\abstract}{\null\vfil}{}{}{} %Abstract patch: do nog delete
\usepackage{tocloft} %Package to properly format Table of Contents
\setlength{\cftbeforetoctitleskip}{0pt}
\setlength{\cftaftertoctitleskip}{0pt}
\usepackage[a4paper, margin=2.54cm]{geometry} %Specifies the size of text area on the paper
\usepackage{hyperref} %Package to make the Table of Contents use hyperlinks
\usepackage{etoolbox} %Package to remove abstract indent
\ifundef{\abstract}{}{\patchcmd{\abstract}%
    {\quotation}{\quotation\noindent\ignorespaces}{}{}}
\usepackage{lipsum} %Package to generate dummy text   
\usepackage{float} %Package to fix figure placement

%\\------------------------------------------------------\\


%\\------------------------------------------------------\\
%\\Add user packages here\\
%\usepackage{natbib} %Package to use Harvard style bibliography, comment out if using IEEE style
%\usepackage{showframe} %Helper package to show document outlines
\usepackage{enumitem} %Package to change the seperation of items in a list

\usepackage{nomencl}%Package to create a list of acronyms
\makenomenclature %Instruction for LaTeX to compile the nomenclature
\renewcommand{\nomname}{\Large List of Abbreviations} %Renames the ToC heading
%\\------------------------------------------------------\\


%\\------------------------------------------------------\\
%\\ Font Specifier
%\\ Comment out this section if you want standard Times New Roman Font
\usepackage{helvet} %Arial Font
\renewcommand{\familydefault}{\sfdefault} %Command needed to make font arial
%\\------------------------------------------------------\\


%\\------------------------------------------------------\\
%\\ Formatting of headings
%\\ Here you can specify the size and color of all headings
\sectionfont{\LARGE }
\subsectionfont{\Large }
\subsubsectionfont{\large }
\renewcommand\abstractname{\LARGE Abstract} %To make the abstract heading purple
\renewcommand{\contentsname}{\center \Large Table of Contents} 
\renewcommand{\listfigurename}{\center\ \Large List of Figures}
\renewcommand{\listtablename}{\center \Large List of Tables }      
\renewcommand\thesection{\arabic{section}} %Specifies the numbering of sections
\renewcommand{\refname}{Bibliography} %Change this to your preferred heading (i.e. Bibliography, List of References etc)
%\\------------------------------------------------------\\


%\\------------------------------------------------------\\
%\\ Page formatting
\setlength{\parindent}{0in} 
\setlength{\parskip}{8pt} %The space between paragraphs
\renewcommand{\baselinestretch}{1.08} %The line spacing
\graphicspath{ {cover/} } %Directory for cover imag
%\\------------------------------------------------------\\


%\\------------------------------------------------------\\
%\\ NWU RGB Purple Color definition
%\\ DO NOT CHANGE THIS
\definecolor{Puk_Pers}{RGB}{108, 61, 145}
%\\------------------------------------------------------\\


%\\------------------------------------------------------\\
%\\ Reformatting of Abstract
\let\oldabstract\abstract
\let\oldendabstract\endabstract
\makeatletter
\renewenvironment{abstract}
{\renewenvironment{quotation}%
               {\list{}{\addtolength{\leftmargin}{-9mm} % change this value to add or remove length to the the default
                        \listparindent 0em%
                        \itemindent    \listparindent%
                        \rightmargin   \leftmargin%
                        \parsep        \z@ \@plus\p@}%
                \item\relax}%
               {\endlist}%
\oldabstract}
{\oldendabstract}
\makeatother
%\\------------------------------------------------------\\


%\\------------------------------------------------------\\
%\\ Nomenclature setup.
%\\ This creates categories for the nomenclature (i.e. Maths, Abbrev etc)
%\\ In order for nomenclature to work properly the following command must be run:
%\\ pdflatex -synctex=1 -interaction=nonstopmode %.tex|pdflatex -synctex=1 -interaction=nonstopmode %.tex|pdflatex -synctex=1 -interaction=nonstopmode %.tex|makeindex %.nlo -s nomencl.ist -o %.nls|pdflatex -synctex=1 -interaction=nonstopmode %.tex|"C:/Program Files/Adobe/Reader 11.0/Reader/AcroRd32.exe" %.pdf
%\\ Add this command to the list of custom user commands under the quick build section in TexMaker
\usepackage{xstring}
\usepackage{xstring}
\usepackage{xpatch}
\patchcmd{\thenomenclature}
  {\leftmargin\labelwidth}
  {\leftmargin\labelwidth\itemindent 1em }
  {}{}
 
\newcommand{\nomenclheader}[1]{%
  \item[\hspace*{-\itemindent}\normalfont\bfseries#1]}
\renewcommand\nomgroup[1]{%
  \IfStrEqCase{#1}{%
  %Here you can change the nomenclature update
   {V}{\nomenclheader{Mathematical terms}}%   V - Variables
   {P}{\nomenclheader{Parameters}}% P - Parameters
   {A}{\nomenclheader{Acronyms}}% A - Acronyms
  }%
}
%\\------------------------------------------------------\\

\begin{document}

\begin{titlepage}

%\\------------------------------------------------------\\
%\\This is for the purple splash
\begin{tikzpicture}[remember picture, overlay, transform shape]
\node [anchor=north west, inner sep=0pt]
    at (current page.north west)
    {\includegraphics[width=\paperwidth]{NWU_Purple_Cover.png}};
\end{tikzpicture}	
%\\------------------------------------------------------\\


%\\------------------------------------------------------\\
%\\ Cover material content
\center
\vspace{4cm}

{\Huge PROJECT TITLE\par}

\vspace{2cm}
	
{\Large\itshape INITIALS AND SURNAME\par}
	
\vspace{2cm}
	
{\LARGE DOCUMENT TYPE HERE  e.g. Practical Report\par}
		
\vspace{5cm}
	
\begin{flushleft}

\large	Supervisor: SUPERVISOR NAME\par
\large	Date: \today  \par %Do not remove the \today command. This automatically changes the date on the cover page to the current date
\large	Student No: STUDENT NO\par
\end{flushleft}

\end{titlepage}
%\\------------------------------------------------------\\

%\\------------------------------------------------------\\
%\\Abstract Section
\thispagestyle{plain} 
\pagenumbering{roman} 
\begin{abstract}
\addcontentsline{toc}{section}{Abstract}
\justify

\lipsum[1] %generates example text


\end{abstract}


\newpage %Use newpage here instead of \pagebreak

\tableofcontents

\listoffigures

\listoftables

\printnomenclature

\pagebreak

\pagenumbering{arabic}

%\\------------------------------------------------------\\
%\\Content Section
\section{Introduction}
\lipsum[1] %generates example text

\subsection{Subsection 1}
\lipsum[1] %generates example text

\subsubsection{SubSubsection 1}
\lipsum[1] %generates example text

\nomenclature[A]{DC}{Direct Current} %Use this to add an entry into the nmeclature section. Refer to the preamble section that deals with the setup of the nomenclature for the index specifier. The letter in square brackets is for grouping of the abbreviations into seperate categories i.e. maths constants and the abbreviations etc

%\cite{BIBLIOGRAPHY_ITEM} - Citation example %%Do NOT use citep with IEEE, it generates compilation errors.

%\\ Uncomment this section for bibliography however make sure that the filename in \bibliography{} matches the actual filename of the .bib file
%\\Use Zotero with the Better BibLaTeX add-on installed for the easiest bibliography setup
%\addcontentsline{toc}{section}{Bibliography} %References/Bibliography has to be manually added to ToC
%\bibliographystyle{IEEEtran} %- IEEE Style
%\bibliographystyle{humannat} %- Harvard style
%\bibliography{BIB FILE NAME}

%\\------------------------------------------------------\\
%\\Common documents extras and how to format them in LaTeX

%\\Figure
%\begin{figure}[H] %It is recommended to use the H specifier. This forces LaTeX to place the figure here and it locks it in place. I.e. No Microsoft Word page formatting dubstep
%\center
%\includegraphics[scale=0.5]{NWU_Purple_Cover.png}
%\caption{CAPTION}
%\label{LABEL_NAME} %LABEL MUST ALWAYS BE AFTER CAPTION AWLAYS
%\end{figure}

%With reference to figure \ref{LABEL_NAME} %This is how you reference a figure/table/section/subsection etc Anything that has a \label is referred to in this way

%Numbered Equation 
%This is to insert a numbered equation
%\begin{equation} 
%x^2 + 4x + 5 
%\label{EQUATION NAME}
%\end{equation}

%Inline Mathematics
%$ x + 1 = 12 $ %use the $ sign for maths you want in your paragraphs

%Tables
%For a detailed guide to tables use https://www.overleaf.com/learn/latex/Tables
%To generate tables easily make use of EITHER the built in table wizard in Texmaker or https://tablesgenerator.com/
%If using the generator (from online or Texmaker) above remember to add caption, label and \center to the table code just underneath the \begin{table}
%e.g.
%\begin{table}[h!]
%\center
%\caption{CAPTION}
%\vspace{0.2cm}
%\label{LABEL NAME}
%\begin{tabular}{|l|l|}
%\hline
%a & b \\ \hline
%c & d   \\ \hline
%\end{tabular}
%\end{table}



%\\------------------------------------------------------\\

\end{document}
